\chapter{Background}

\section{Vehicular communications and spectrum reuse}
% Lo necesario de comunicaciones inalámbricas para entender el modelo: interferencia, calidad de una señal y relación entre cola y retardo.

\subsection{Vehicle-to-vehicle links and the cellular network}
% Qué son V2X y V2V, qué es un par V2V, qué mensajes lleva y cómo reutiliza los canales de subida de los CUE a costa de interferirlos.

\subsection{Power and decibels}
% Escala logarítmica: razones de potencia en dB y potencias absolutas en dBm, con ejemplos de conversión.

\subsection{Channel gain: distance and fading}
% La ganancia h = alpha * g: pérdida por distancia (gran escala) y desvanecimiento rápido (pequeña escala).

\subsection{Signal-to-interference-plus-noise ratio}
% Definición de la SINR, potencia del ruido térmico e interferencia de otros transmisores en el mismo canal.

\subsection{Channel capacity}
% Fórmula de Shannon, R = W log2(1 + SINR), y cómo crece la tasa con la SINR.

\subsection{Queues, delay, and Little's law}
% Cómo crece y se vacía una cola según la tasa de llegada y la de transmisión, y la ley de Little, que relaciona el tamaño de la cola con el retardo.

\section{Reinforcement learning}
% Cómo aprende un agente por refuerzo, desde la interacción con el entorno hasta Q-learning.

\subsection{Agent, environment, and reward}
% El ciclo estado, acción, recompensa y nuevo estado, y qué es una política.

\subsection{Return and discount factor}
% El retorno como suma descontada de recompensas y el papel del factor de descuento gamma.

\subsection{Markov decision processes}
% Los cinco elementos de un MDP y la propiedad de Markov.

\subsection{Value functions and the Bellman equation}
% La función Q de una política, la función Q óptima y la ecuación de Bellman.

\subsection{Q-learning}
% La regla de actualización de Q-learning y el error de diferencia temporal (TD).

\subsection{Exploration versus exploitation}
% La política epsilon-greedy y cómo se reduce epsilon durante el entrenamiento.

\section{Deep reinforcement learning}
% Cómo se combina el aprendizaje por refuerzo con redes neuronales y las mejoras que forman el agente D3QN.

\subsection{Neural networks as function approximators}
% Por qué no sirve una tabla cuando el estado es continuo y cómo una red neuronal con pesos theta aproxima la función Q.

\subsection{Deep Q-networks}
% La pérdida de la DQN, el buffer de repetición de experiencia y la red objetivo Mnih et al.

\subsection{Double DQN}
% Por qué el máximo sobrestima los valores y cómo la doble DQN separa elegir la acción de evaluarla.

\subsection{Dueling architecture}
% Estimar por separado el valor del estado V(s) y la ventaja A(s, a), y cómo se combinan en Q.

\subsection{Prioritized experience replay}
% Muestrear las transiciones según su error TD, el exponente de prioridad y los pesos de importancia que corrigen el sesgo.

\subsection{Huber loss}
% Pérdida cuadrática para errores pequeños y lineal para errores grandes, para evitar saltos bruscos en los pesos.

\subsection{The D3QN network}
% Cómo se combinan duelo, doble DQN y PER en el agente que usa cada par V2V.

\section{Multi-agent reinforcement learning}
% Qué cambia cuando aprenden varios agentes a la vez en el mismo entorno.

\subsection{From one agent to many}
% Agentes que comparten un entorno, la acción conjunta y la formalización como Dec-POMDP.

\subsection{Why it is hard}
% Las cuatro dificultades propias del caso multiagente: observabilidad parcial, no estacionariedad, escalabilidad y asignación de crédito.

\subsection{Ways to train multiple agents}
% Entrenamiento y ejecución centralizados, CTDE y DTE (clasificación de Amato), y dónde encaja el aprendizaje federado.

\section{Federated learning}
% La idea del aprendizaje federado y su aplicación al aprendizaje por refuerzo.

\subsection{The idea and the FedAvg algorithm}
% Entrenar sin mover los datos: copias locales del modelo, un servidor que promedia y la fórmula de FedAvg (McMahan et al).

\subsection{Federated reinforcement learning}
% Aplicar el promedio federado a agentes de refuerzo; FRL horizontal, que es el caso que uso.

\subsection{The cost of the central server}
% El servidor como punto único de falla, cuello de botella de la comunicación y entidad en la que todos deben confiar.

\section{Decentralized federated learning and gossip}
% Cómo aprender sin servidor, intercambiando modelos solo entre vecinos.

\subsection{Communication graph}
% Nodos, aristas, vecinos y grado; grafos geométricos aleatorios como modelo de una red inalámbrica.

\subsection{Averaging with neighbors}
% Algoritmos de gossip: una ronda de promedio con los vecinos y cómo la estructura del grafo determina la velocidad de convergencia al consenso.

\subsection{Learning and averaging at the same time}
% Combinar pasos de entrenamiento local con rondas de promedio (D-PSGD de Lian et al) y su costo de comunicación frente a la versión centralizada.
